
\documentclass[11pt]{article}

\usepackage{natbib,hyperref,ragged2e,graphicx}


% ref
% 


\title{Denise Pumain, visions and innovations pushing forward Open Science}
\author{J. Raimbault$^{1,2,3,4,\ast}$\medskip\\
$^{1}$ LaSTIG, Universit{\'e} Gustave Eiffel, IGN-ENSG\\
$^{2}$ CASA, UCL, United Kingdom\\
$^{3}$ UPS CNRS 3611 ISC-PIF, France\\
$^{4}$ UMR CNRS 8504 G{\'e}ographie-cit{\'e}s, France\\
$\ast$ \texttt{juste.raimbault@ign.fr}
}

\date{}


\begin{document}


{\justify \textit{Published as: } Raimbault, J. (2023). Denise Pumain, visions and innovations pushing forward Open Science. Cybergeo: European Journal of Geography. \url{https://journals.openedition.org/cybergeo/40196}}


{\let\newpage\relax\maketitle}

\begin{abstract}
	
\end{abstract}


\bigskip


I had the privilege to start working with Denise Pumain for the first time in late 2015, as my PhD supervisor Arnaud Banos had suggested me as a contributor to the ``Cybergeo 20'' initiative. She was leading an innovative data analysis initiative, aiming at enhancing reflexivity for readers and contributors to the journal. To celebrate the 20th anniversary of the Cybergeo journal, Denise Pumain had the vision to construct corpus analysis tools that would be applied to the content of the journal themselves. Besides the prototype web application available online (at \url{https://analytics.huma-num.fr/geographie-cites/cybergeonetworks/}) and combining the diverse analysis led by the interdisciplinary team from the Géographie-cités laboratory (Pierre-Olivier Chasset, Clémentine Cottineau, Hadrien Commenges, and myself), the initiative was synthesised into several publications  \citep{cottineau2017cybergeonetworks,raimbault2021empowering,raimbault2019exploration}, highlighting the timeliness of Denise’s idea. Indeed, in the following years, similar tools (related themes or papers for example) started to appear on commercial publishers websites. When metadata become the new eldorado for the interests of publishers, keeping the power on such reflexive analysis is crucial for the survival of truly open initiatives such as Cybergeo.


Of course, Denise’s contributions to advocate and act for Open Science go back far more in time, starting with the funding of Cybergeo itself in 1996, which was the first fully electronic geography journal. In 2002, while social science and geography was lagging in the transition towards online and open journals, Cybergeo was still the only geography online journal, fully open and with scientific standards of publication \citep{kosmopoulos2002cybergeo}. Cybergeo has always been able to change and adapt, updating its organisation to enhance interactivity and sollicitate debates \citep{pumain1999cybergeo}. More recently in 2019, the CybergeoNet project lead by Christine Kosmopoulos, aimed at facilitating the translation of metadata and some papers, through a workflow transposable to similar journals. Open Access journals indeed play a central role in the international diffusion of scientific knowledge \citep{pumain2022role}. The Hypergeo online open encyclopedia (\url{https://hypergeo.eu/}) for geography was also involved in this project, and it is no surprise that Denise is involved in this open project, as a core funding member, but also as the most prolific author in the encyclopedia (121 articles out of around 500 in total).



Also under the impulsion of Denise, Cybergeo has constantly adapted to follow the best practices in Open Science. Data citation is an important aspect of the F.A.I.R. principles – which through Findability, Accessibility, Interoperability and Reuse sets guidelines for data management in an Open Science framework (Parsons et al., 2019). Cybergeo has thus opened in 2017 a specific section ``Data Papers'', with the aim to encourage data sharing, publication and citation for geography \citep{pumain2017expressions}. While models are also a core component of scientific knowledge production, at least in quantitative approaches, their description and citation usually remains linked to an complete paper, which in some cases may hinder the preliminary debates taking place during model construction, or the flexibility and reusability of the model. Therefore, the ``Model Papers'' section of Cybergeo (GeoOpenMod), pushed forward by Arnaud Banos, aims at facilitating model sharing, reuse, and validation, but also at providing a basis for citation for open source software in geography \citep{banos2014introduction}.



Denise Pumain’s contributions to a more open and robust science go beyond the domain of science communication, with a deep impact on modelling and simulation practices. The ERC project GeoDivercity, of which Denise was the Principal Investigator between 2010 and 2015, significantly advanced knowledge on worldwide systems of cities, but also introduced new methods and tools to explore and validate simulation models \citep{pumain2017urban,pumain2020theories}. The urban dynamics agent-based models were put forward as open models, such as the Marius modular framework \cite{cottineau2015modular}. Furthermore, the open source and free platform OpenMOLE \citep{reuillon2013openmole}, although existing since 2008, has been significantly developed through the project and the collaboration with geographers. This software for model validation is based on three complementary axis: model embedding with an agnosticity to programming langage, transparent access to high performance computing infrastructure, and state-of-the-art methods in sensitivity analysis and model validation. First, it allows thematic researcher to access such resources and methods, thus opening the scope of model validation practices. Second, model benchmark and validation, allowed by the methods provided, is an important step before being able to share and reuse models. Third, the scripting language makes it easy to reproduce computational experiments but also to couple models. Such practices are a core component for Open Science, and OpenMOLE has enabled them in social science and geography. A present example illustrates well this link between Open Science and model exploration. The French national mapping Institute IGN is currently elaborating a project of Digital Twin (\url{https://www.ign.fr/les-communs-dutilite-publique}), in the context of geocommons (shared and open geographical data). For the simulation part, OpenMOLE is considered to embed, couple, link with data, explore and validate simulation models that will be integrated into the Twin. Thus, these open simulation practices which are rather new in our fields, have been largely put forward and advanced thanks to Denise’s contributions.



In terms of open data construction and diffusion, Denise Pumain’s contributions must also be recalled. For example, the construction of the first open database for Chinese cities population with a consistent geographical definition, was achieved with her PhD student Elfie Swerts \citep{swerts2010peut,swerts2017data}. The database for large urban systems worldwide, constructed during the GeoDivercity project, and which in particular ensure comparability between the systems, is also an open database \citep{pumain2015multilevel,cura2017old}, available through the GibratSim application \citep{cura2017gibratsim}.



Finally, current ongoing developments and future perspectives, related to aforementioned projects, originate with Denise’s recent proposals. The cybergeonetworks software developed through the ``Cybergeo 20'' project remains a prototype and can currently not scale, automatically update, or be transposed to other journals. A project starting in 2023 will aim at integrating its functionalities into a plugin for the Gargantext software (\url{https://gargantext.org/}). This open and free software is a robust and modular tool for text-mining, what will solve scalability and portability problems. This should provide an open and flexible tool that can be used by independent journal editors to provide their readership with literature mapping and corpus mining possibilities to explore the scientific environment of the journal, crucial for reflexivity in an open science context. An other perspective for Cybergeo currently discussed and which stems from discussions with Denise, is to use the model paper section to foster reproducibility and validity of exchanged models. In its current form, only source code deposited on an open repository and a detailed description of the model are required. This does not ensure complete reproducibility, as many problems with software environments, versions, or even missing data, may occur when trying to run the model. The use of docker container is a solution to ensure reproducibility \citep{boettiger2015introduction}, and may be an additional requirement to ask to authors. More generally, best practices for computational modelling \citep{stodden2014best} may be required by this journal section in a near future. Furthermore, pure replicability of numerical results is not enough for reproducibility, and qualitative results should be confirmed by independent numerical experiments with possibly different designs of experiments. Thus, the efforts of the author to validate the model should be extensively detailed in a model paper, to facilitate future experiments for further validation. In the case where OpenMOLE is used, scripts can be shared along the model. Of course, there is no universal definition of validation across disciplines and contexts (and a broad and systematic interdisciplinary view on this concepts still lacks in the literature), such that no particular method or experiment should be imposed on authors. On the contrary, they should be co-constructed with the model itself \citep{chapron2022co}. There is therefore still a long way to go before finding the correct compromise to foster model validation practices without being too normative. The Cybergeo journal will surely be at the forefront of such developments and experiments, within the context of a more general activism for Open Science, and without a doubt owing a lot to Denise’s legacy after she will retire as journal chief editor in February 2023.


I would like to thank Denise Pumain for all these efforts, innovative ideas, and concrete actions to contribute to the mutation of geography into an open science, and more generally to the advent of Open Science itself. This echoes with the slogan of the open software community Framasoft, ``\textit{La route est longue mais la voie est libre}'', but to which I would like to add: thanks to visionaries like Denise, we know in which direction this road should be built.


\bigskip


\bibliographystyle{apalike}
\bibliography{biblio}

%\begin{thebibliography}{}
%
%\end{thebibliography}



\end{document}

